\documentclass[10pt,a4paper]{amsart}
\usepackage[margin=19mm]{geometry}
\usepackage{amsmath,amssymb,mathtools}
\usepackage[scr=rsfs,cal=euler]{mathalfa}
\usepackage{xfrac}
\usepackage[unicode,hidelinks,bookmarks=false]{hyperref}
\newcommand{\ie}{i.e.\ }
\newcommand{\cf}{cf.\ }
\newcommand{\resp}{respectively\ }
\newcommand{\Subsection}{\subsection}
\providecommand{\nobreakdash}{\nobreak\hbox{-}}
\setlength{\emergencystretch}{2em}
\multlinegap0pt
\usepackage{bm}
\usepackage{bbm}
\usepackage{dsfont}
\usepackage{xspace}
\usepackage{cancel}
\usepackage{mathtools}
\usepackage{amsthm}
\makeatletter
\@ifundefined{assumption}{\newtheorem{assumption}{Assumption}}{}
\@ifundefined{lemma}{\newtheorem{lemma}{Lemma}}{}
\@ifundefined{proposition}{\newtheorem{proposition}{Proposition}}{}
\makeatother
\title[Asymptotic Theory for Linear Functionals of Kernel Ridge Reg]{Asymptotic Theory for Linear Functionals of Kernel Ridge Regression}
\author{Rui Tuo, Lu Zou}
\date{Source: arXiv:2403.04248v3 (CC BY 4.0)}
\begin{document}
\maketitle

\noindent\emph{Source and licence.} Asymptotic Theory for Linear Functionals of Kernel Ridge Regression. Version arXiv:2403.04248v3 (CC BY 4.0), distributed under the Creative Commons Attribution 4.0 International licence, as recorded in the arXiv licence metadata for this exact version and shown on the abstract page https://arxiv.org/abs/2403.04248v3. Full rights notice: licenses/arxiv-2403.04248-CC-BY-4.0.txt.\par\smallskip

\noindent\textit{Editorial scope.} Counted unit: the Lemma whose source label is \texttt{lemma:step2} in the article (source0.tex lines 1463--1475, source label \texttt{lemma:step2}), delivered together with the article's own complete original proof of it (source0.tex lines 1477--1503, one \texttt{proof} environment of 27 lines). No mathematics is written, repaired or re-derived: statement and proof are the article's own text line for line. Required context retained uncounted: assum:matern (lines 373--376); interpolation (lines 566--568); assum:random1 (lines 1386--1388); lemma:bernstein (lines 1423--1426); entropy (lines 1454--1460). Every definition, display and label of those lines is retained unchanged. Omitted from this package and not redistributed: 1--372, 377--565, 569--1385, 1389--1422, 1427--1453, 1461--1462, 1476--1476, 1504--2355. Nothing inside a retained excerpt is omitted or altered: every retained body is delivered exactly as the article's lines are written, so no omission receipt of any kind accompanies this package. Citations: the retained excerpts carry no citation command at all, so no bibliography entry of the article is transcribed and no reference is redistributed. Reference adaptation: none was needed and none was made; every reference inside the delivered unit resolves inside it, so no unresolved reference or citation is produced. Printed slips kept literal: no printed word, symbol, hypothesis or number of the counted statement or of its proof was altered. Layout and class: the article's own class is \texttt{article} with packages including amsthm, amsmath, amsfonts, amssymb, mathrsfs, natbib, graphicx, adjustbox, float, caption, subcaption, booktabs, makecell, multirow; the wrapper is the lane's pinned \texttt{amsart} editorial wrapper at 10pt on A4, which restates the theorem-like environments the retained text needs and adds the article's own macro definitions that the retained text needs (none), with extra wrapper packages (bm, bbm, dsfont, xspace, cancel, mathtools). The delivered numbering is the unit's own. Assets: no retained span contains a figure environment, a graphics-inclusion command, a PDF-page inclusion, a table or a TikZ picture; no image file is copied into this package, the package declares no asset dependency and no image file is redistributed with it. Licence: version arXiv:2403.04248v3 of this article is distributed under the Creative Commons Attribution 4.0 International licence, as recorded in the arXiv licence metadata for this exact version and shown on the abstract page https://arxiv.org/abs/2403.04248v3. A full rights notice is delivered at licenses/arxiv-2403.04248-CC-BY-4.0.txt. Characters: every retained excerpt is pure ASCII, so the delivered engine, XeLaTeX with its default Latin Modern text font, reports no missing character.
\begin{assumption}\label{assum:matern}
    The input domain $\Omega$ is a convex and compact subset of $\mathbb{R}^d$ with a non-empty interior.
    In addition, $\mathcal{H}$ is equal to a (fractional) Sobolev space with order $m$ (satisfying $m>d/2$), denoted by $H^m$, with equivalent norms. 
\end{assumption}

\begin{eqnarray}\label{interpolation}
    \|v\|_{L_\infty}\leq A\|v\|^{1-\frac{d}{2m}}_{L_2}\|v\|^{\frac{d}{2m}}_{H^m},
\end{eqnarray}

\begin{assumption}\label{assum:random1}
    The input sites $x_1,\ldots,x_n$ are independent and identically distributed random variables over $\Omega$ with density function $\mu(\cdot)$. In addition, $\inf_{x\in\Omega}\mu(x)=\mu_0>0$.
\end{assumption}

\begin{lemma}\label{lemma:bernstein}
    Suppose $v$ is continuous over $\Omega$. Under Assumption \ref{assum:random1}, we have
    \[\mathbb{P}(\|v\|_{L_2}\leq \sqrt{2/\mu_0}\|v\|_n)\geq 1- \exp\left\{-\frac{3n\mu_0\|v\|^2_{L_2}}{28\|v\|^2_{L_\infty}}\right\}.\]
\end{lemma}

\begin{proposition}
Suppose $\Omega$ satisfies the conditions in Assumption \ref{assum:matern}.
    There exists a constant $A>0$ depending only on $\Omega,m,d$, such that all $r>0$,
    \begin{eqnarray}\label{entropy}
        \log N(r,H^m(R))\leq A(R/r)^{d/m}.
    \end{eqnarray}
\end{proposition}

\begin{lemma}\label{lemma:step2}
Suppose $\Omega$ satisfies the conditions in Assumption \ref{assum:matern}.
    Fix $R>0$. For any $r>0$ satisfying 
    \begin{eqnarray}
        \sqrt{n}(r/R)^{d/m}\geq 1,\label{epsilon}
    \end{eqnarray}
    under Assumption \ref{assum:random1},
   there exists constants $C_1,C_2,C_3,C_4$ depending only on $\Omega,m,d$ and $\mu_0$, such that
   \begin{eqnarray*}
       \mathbb{P}\left(\|v\|_{L_2}\leq C_1 \|v\|_n \text{ for all } v\in H^m(R) \text{ with } \|v\|_{L_2}\geq C_2r\right)\\\geq 1-C_3\exp\{-C_4n(r/R)^{d/m}\}
   \end{eqnarray*}
  % \[\mathbb{P}\left(\|v\|_{L_2}\leq C_1 \|v\|_n \text{ for all } v\in H^m(R) \text{ with } \|v\|_{L_2}\geq C_2r\right)\geq 1-C_3\exp\{-C_4n(r/R)^{d/m}\}.\]
\end{lemma}

\begin{proof}
    The proof proceeds by applying a peeling device.
    Let $|\Omega|$ be the volume of $\Omega$, and $\mathcal{V}_s:=\{v\in\mathcal{V}(R):(s-1)|\Omega|^{1/2}r \leq \|v\|_{L_2}\leq s |\Omega|^{1/2}r\}$ for $s=1,2,\ldots$. By the definition of the covering number, we have $N(r,\mathcal{V})$ centers. For $v\in\mathcal{V}_s$, denote its associated center as $\operatorname{ctr}v$. Note that
    \[(s-2)|\Omega|^{1/2}r\leq\|v\|_{L_2}-|\Omega|^{1/2}r\leq \|\operatorname{ctr}v\|_{L_2}\leq\|v\|_{L_2}+|\Omega|^{1/2}r\leq (s+1)|\Omega|^{1/2}r.\] Define event
    \[E_v:=\{\|\operatorname{ctr}v\|_{L_2}\leq \sqrt{2/\mu_0}\|\operatorname{ctr}v\|_n\}.\]
    Then on $E_v$, we have 
    \begin{eqnarray*}
        \|v\|_{L_2}&\leq& s|\Omega|^{1/2}r\leq \frac{s}{s-2}\|\operatorname{ctr}v\|_{L_2}\leq \frac{s\sqrt{2/\mu}}{s-2}\|\operatorname{ctr}v\|_n\\
        &\leq&\frac{s\sqrt{2/\mu_0}}{s-2}(\|v\|_n+r)\leq \frac{s\sqrt{2/\mu_0}}{s-2}\left(\|v\|_n+\frac{|\Omega|^{-1/2}}{s-1}\|v\|_{L_2}\right).
    \end{eqnarray*}
    Then for $s> 2|\Omega|^{-1/2}+1$, we have $\|v\|_{L_2}\leq 4\sqrt{2/\mu_0}\|v\|_n$. This proves $E_v\subset \{\|v\|_{L_2}\leq 4\sqrt{2/\mu_0}\|v\|_n\}$. Therefore, by Lemma \ref{lemma:bernstein}, we have
    \begin{eqnarray*}
    &&\mathbb{P}(\|v\|_{L_2}> 4\sqrt{2/\mu_0}\|v\|_n)\leq \mathbb{P}(E_v^c)\leq \exp\left\{-\frac{3n\mu_0\|v\|^2_{L_2}}{28\|v\|^2_{L_\infty}}\right\}\\
    &\leq&\exp\left\{-\frac{3n\mu_0\|v\|^2_{L_2}}{28C\|v\|^{2-d/m}_{L_2}\|v\|^{d/m}_{H^m}}\right\}=\exp\left\{-\frac{3n\mu_0}{28C}\frac{\|v\|^{d/m}_{L_2}}{\|v\|^{d/m}_{H^m}}\right\}\\
    &\leq&\exp\left\{-\frac{3n\mu_0}{28C}\left(\frac{(s-1)|\Omega|^{1/2}r}{R}\right)^{d/m}\right\},
    \end{eqnarray*}
    where the third inequality follows from the interpolation inequality (\ref{interpolation}). Choose $S_0$ large enough such that we also have $3\mu_0(S_0-1)^{d/m}|\Omega|^{d/(2m)}/(28C)>(A+1)S_0^{d/(2m)}$, for $A$ defined in (\ref{entropy}). Now we arrive at
    \begin{eqnarray*}
        &&\mathbb{P}\left(\bigcup_{v\in \cup_{s\geq S_0}\mathcal{V}_s}\left\{\|v\|_{L_2}> 4\sqrt{2/\mu_0}\|v\|_n\right\}\right)\\
        &\leq&\sum_{s\geq S_0}\mathbb{P}\left(\bigcup_{v\in \mathcal{V}_s}\left\{\|v\|_{L_2}> 4\sqrt{2/\mu_0}\|v\|_n\right\}\right)\\
        &\leq&\sum_{s\geq S_0}\mathbb{P}\left(\bigcup_{v\in \mathcal{V}_s}E_v^c\right)\\
        &\leq&\sum_{s\geq S_0}\exp\left\{\log N(r,H^m(R))-\frac{3n\mu_0}{28C}\left(\frac{(s-1)|\Omega|^{1/2}r}{R}\right)^{d/m}\right\}\\
        &\leq&\sum_{s\geq S_0}\exp\left\{A(R/r)^{d/m}-(A+1)s^{d/(2m)}n(r/R)^{d/m}\right\}\\
        &\leq&\sum_{s\geq S_0}\exp\left\{-s^{d/(2m)}n(r/R)^{d/m}\right\},
    \end{eqnarray*}
    where the last inequality follows from (\ref{epsilon}). This completes the proof.
\end{proof}

\end{document}
